\begin{thebibliography}{99}
\footnotesize
\setlength{\itemsep}{0.1em}

\bibitem{hong2021} 洪永淼, 汪寿阳. 大数据如何改变经济学研究范式?[J]. 管理世界, 2021, 37(10): 40--55.
\bibitem{suzhi2017} 苏治, 卢曼, 李德轩. 深度学习的金融实证应用: 动态、贡献与展望[J]. 金融研究, 2017(5): 111--126.
\bibitem{wangfang2020} 王芳, 王宣艺, 陈硕. 经济学研究中的机器学习: 回顾与展望[J]. 数量经济技术经济研究, 2020, 37(4): 146--164.
\bibitem{jiang2017} Jiang Z, Xu D, Liang J. A deep reinforcement learning framework for the financial portfolio management problem[EB/OL]. arXiv:1706.10059, 2017.
\bibitem{liu2020} Liu X Y, Yang H, Chen Q, et al. FinRL: A deep reinforcement learning library for automated stock trading in quantitative finance[EB/OL]. arXiv:2011.09607, 2020.
\bibitem{li2024} Li Z, Tam V, Yeung K. Developing a multi-agent and self-adaptive framework with deep reinforcement learning for dynamic portfolio risk management[EB/OL]. arXiv:2402.00515, 2024.
\bibitem{grossman1980} Grossman S J, Stiglitz J E. On the impossibility of informationally efficient markets[J]. American Economic Review, 1980, 70(3): 393--408.
\bibitem{mnih2015} Mnih V, Kavukcuoglu K, Silver D, et al. Human-level control through deep reinforcement learning[J]. Nature, 2015, 518(7540): 529--533.
\bibitem{lillicrap2016} Lillicrap T P, Hunt J J, Pritzel A, et al. Continuous control with deep reinforcement learning[C]//International Conference on Learning Representations, 2016.
\bibitem{sutton2018} Sutton R S, Barto A G. Reinforcement Learning: An Introduction[M]. 2nd ed. Cambridge: MIT Press, 2018.
\bibitem{charpentier2023} Charpentier A, Élie R, Remlinger C. Reinforcement learning in economics and finance[J]. Computational Economics, 2023, 62(1): 425--462.
\bibitem{cong2021} Cong L W, Tang K, Wang J, et al. AlphaPortfolio: Direct construction through deep reinforcement learning and interpretable AI[EB/OL]. 2021. SSRN:3554486.
\bibitem{hamilton1989} Hamilton J D. A new approach to the economic analysis of nonstationary time series and the business cycle[J]. Econometrica, 1989, 57(2): 357--384.
\bibitem{guidolin2007} Guidolin M, Timmermann A. Asset allocation under multivariate regime switching[J]. Journal of Economic Dynamics and Control, 2007, 31(11): 3503--3544.
\bibitem{timmermann2006} Timmermann A. Forecast combinations[M]//Elliott G, Granger C W J, Timmermann A. Handbook of Economic Forecasting, Vol. 1. Amsterdam: Elsevier, 2006: 135--196.
\bibitem{ledoit2004} Ledoit O, Wolf M. Honey, I shrunk the sample covariance matrix[J]. Journal of Portfolio Management, 2004, 30(4): 110--119.
\bibitem{demiguel2009} DeMiguel V, Garlappi L, Nogales F J, et al. A generalized approach to portfolio optimization: Improving performance by constraining portfolio norms[J]. Management Science, 2009, 55(5): 798--812.
\bibitem{sims2003} Sims C A. Implications of rational inattention[J]. Journal of Monetary Economics, 2003, 50(3): 665--690.
\bibitem{dellavigna2009} DellaVigna S, Pollet J M. Investor inattention and Friday earnings announcements[J]. Journal of Finance, 2009, 64(2): 709--749.
\bibitem{odean1998} Odean T. Are investors reluctant to realize their losses?[J]. Journal of Finance, 1998, 53(5): 1775--1798.
\bibitem{barberodean2001} Barber B M, Odean T. Boys will be boys: Gender, overconfidence, and common stock investment[J]. Quarterly Journal of Economics, 2001, 116(1): 261--292.
\bibitem{hongstein1999} Hong H, Stein J C. A unified theory of underreaction, momentum trading, and overreaction in asset markets[J]. Journal of Finance, 1999, 54(6): 2143--2184.
\bibitem{shleifer1997} Shleifer A, Vishny R W. The limits of arbitrage[J]. Journal of Finance, 1997, 52(1): 35--55.
\bibitem{gu2020} Gu S, Kelly B, Xiu D. Empirical asset pricing via machine learning[J]. Review of Financial Studies, 2020, 33(5): 2223--2273.
\bibitem{gu2021} Gu S, Kelly B, Xiu D. Autoencoder asset pricing models[J]. Journal of Econometrics, 2021, 222(1): 429--450.
\bibitem{chen2024} Chen L, Pelger M, Zhu J. Deep learning in asset pricing[J]. Management Science, 2024, 70(2): 714--750.
\bibitem{harvey2016} Harvey C R, Liu Y, Zhu H. ... and the cross-section of expected returns[J]. Review of Financial Studies, 2016, 29(1): 5--68.
\bibitem{feng2020} Feng G, Giglio S, Xiu D. Taming the factor zoo: A test of new factors[J]. Journal of Finance, 2020, 75(3): 1327--1370.
\bibitem{kozak2020} Kozak S, Nagel S, Santosh S. Shrinking the cross-section[J]. Journal of Financial Economics, 2020, 135(2): 271--292.
\bibitem{kelly2024} Kelly B T, Malamud S, Zhou K. The virtue of complexity in return prediction[J]. Journal of Finance, 2024, 79(1): 459--503.
\bibitem{kelly2023} Kelly B T, Xiu D. Financial machine learning[J]. Foundations and Trends in Finance, 2023, 13(3--4): 205--363.
\bibitem{libin2019} 李斌, 邵新月, 李玥阳. 机器学习驱动的基本面量化投资研究[J]. 中国工业经济, 2019(8): 61--79.
\bibitem{jiang2021} 姜富伟, 孟令超, 唐国豪. 媒体文本情绪与股票回报预测[J]. 经济学(季刊), 2021, 21(4): 1323--1344.
\bibitem{shen2019} 沈艳, 陈赟, 黄卓. 文本大数据分析在经济学和金融学中的应用: 一个文献综述[J]. 经济学(季刊), 2019, 18(4): 1153--1186.
\bibitem{lopezdeprado2018} López de Prado M. Advances in Financial Machine Learning[M]. Hoboken: John Wiley \& Sons, 2018.
\bibitem{fuster2022} Fuster A, Goldsmith-Pinkham P, Ramadorai T, et al. Predictably unequal? The effects of machine learning on credit markets[J]. Journal of Finance, 2022, 77(1): 5--47.
\bibitem{bommasani2021} Bommasani R, Hudson D A, Adeli E, et al. On the opportunities and risks of foundation models[EB/OL]. 2021. arXiv:2108.07258.
\bibitem{chakraborty2017} Chakraborty C, Joseph A. Machine learning at central banks[R]. Bank of England Staff Working Paper No. 674, 2017.
\bibitem{bis2024aer} Bank for International Settlements. Artificial intelligence and the economy: Implications for central banks[R]//Annual Economic Report 2024, Chapter III. Basel: BIS, 2024.
\bibitem{brunnermeier2009} Brunnermeier M K, Pedersen L H. Market liquidity and funding liquidity[J]. Review of Financial Studies, 2009, 22(6): 2201--2238.

\end{thebibliography}
